\vspace{-0.1in}
\begin{table}[t]
\begin{subtable}[t]{0.3\textwidth}
\small
\begin{tabular}{c}
\toprule[1.5pt]      
Tasks of \textit{LRV-Instruction} \\ 
\midrule
\midrule  
\pbox{7cm}{\textit{Image Captioning, Object Detection, Image Sentiment Analysis, Image Quality Assessment, Object Interaction Analysis, Image Anomaly Detection, Referential Expression Grounding, OCR, VCR, Object Attribute Detection, Muli-choice VQA, Semantic Segmentation, Dense Captioning, Visual Entailment, Styled Image Caption, Activity Recognition, Facial Expression Detection
}}\\
\bottomrule[1.5pt]
\end{tabular}
 \caption{VL tasks in our dataset.}
\label{tab:stat1}
\end{subtable}
\hspace{\fill}
\begin{subtable}[t]{0.5\textwidth}
\setlength\tabcolsep{4.7pt}
\centering
\small
\begin{tabular}{lc}
\toprule[1.5pt]
Statistic &\\
\midrule
\# of Instructions      & 120k\\
\# of Images         & 6.5k\\
\# of Interrogative Instructions        & 70k \\
\# of Declarative Instructions        & 50k \\
Ave \# of Instruction per Image       & 17.7 \\
Ave Answer Length      & 15.2 \\
Ave Instruction Length     & 12.1 \\
\midrule
\# of Negative Instructions     & 35k\\
Ave Instruction Length (Neg1) & 18.1\\
Ave Instruction Length (Neg2) & 11.2\\
\% of Interrogative (Neg1) & 47\% \\
\% of Interrogative (Neg2) & 81\% \\
\bottomrule[1.5pt]
\end{tabular}
\vspace{0.05in}
\caption{
Statistics of \textit{LRV-Instruction}.
}
\label{tab:stat2}
\end{subtable}
\caption{ (a) The proposed \textit{LRV-Instruction} covers a diverse set of vision-language tasks. (b) Detailed statistics of our constructed dataset. 
Neg1 means \textit{Nonexistent Element Manipulation} and Neg2 means \textit{Existent Element Manipulation} in the negative instructions.
}
\label{tab:stat}
\vspace{-0.2in}
\end{table}

CHAIR \cite{rohrbach2018object} was introduced to evaluate object hallucination in image captioning tasks. However, it 
usually demands complex human-crafted rules. Alternatively, \cite{li2023evaluating} formulates the evaluation of object hallucination as a binary classification task that prompts LMM to output "Yes" or "No". However, it is hard to evaluate the LMM output in an open-ended manner. In addition, both methods highly depend on human-annotated groundtruth answers.


To this end, we introduce \textit{GPT4-Assisted Visual Instruction Evaluation (GAVIE)} as a more flexible and robust approach. The general prompt we use is shown in %Fig. \ref{fig:prompt_eval} in 
the Appendix. GPT4 takes the dense captions with bounding box coordinates as the image content and compares human instructions and model response. Then we ask GPT4 to work as a smart teacher and score (0-10) students' answers based on two criteria. \textit{(1) Accuracy: whether the response is accurate concerning the image content}. \textit{(2) Relevancy: whether the response directly follows the instruction}. An evaluation result is shown in Fig. \ref{fig:model}, which successfully points out that \textit{"dog, ball"} is inconsistent with the image, and the response from the first LMM did not address the instruction. Unlike previous evaluation methods \cite{liu2023visual, li2023evaluating, rohrbach2018object}, \textit{GAVIE} does not require human-annotated groundtruth answers and can freely adapt to diverse instruction formats.

\begin{table}[t]
\setlength\tabcolsep{3pt}
\centering
\small
\begin{tabular}{lccccc}
\toprule[1.5pt]
& Ours & MiniGPT4 & LLaVA & InstructBLIP & MMGPT\\
\midrule
\textsc{Accuracy(GPT4)} & \textbf{6.58} & 4.14 & 4.36 & 5.93 & 0.91\\
\textsc{Relevancy(GPT4)} & \textbf{8.46} & 5.81 & 6.11 & 7.34 & 1.79\\
\bottomrule[1.5pt]
\end{tabular}
\vspace{0.05in}
\caption{
Comparison Experiment on our evaluation set evaluated by \textit{GAVIE}. The metric scores of \textsc{Accuracy(GPT4)} and \textsc{Relevancy(GPT4)} are in a scale of 0 to 10.
}
\label{tab:comparison_evaluation}
\vspace{-0.1in}
\end{table}


\begin{table}[t]
\setlength\tabcolsep{3pt}
\centering
\small
\begin{tabular}{lcccc}
\toprule[1.5pt]
Model & Accuracy & F1 Score\\
\midrule
LLaVA & 0.45 & 0.59 \\
% \midrule
MiniGPT4 & 0.44 & 0.52\\
% \midrule
InstructBLIP & 0.74 & 0.81 \\
% \midrule
Ours & \textbf{0.79} & \textbf{0.84} \\
\bottomrule[1.5pt]
\end{tabular}
\vspace{0.05in}
\caption{Comparison experiment on hallucination evaluation dataset introduced by \cite{li2023evaluating}.} 
\label{tab:coco}
\vspace{-0.2in}
\end{table}





